\section{Preliminaries}
\label{sec:prelim}

\subsection{Notation}
Inputs $x$ are drawn from a distribution $D$ on an input space $\mathcal{X}$, and
$f$ denotes a fixed model. Expectations $\E[\cdot]$ are over $x\sim D$ unless
otherwise indicated. We write $\Phi$ for the standard-normal cumulative
distribution function and $\Phi^{-1}$ for its inverse. For $a\le b$ and a real
number $t$ we write $\clip(t,a,b)=\min\{\max\{t,a\},b\}$.

\subsection{Perturbation sets}
An adversary is modeled by a family of \emph{perturbation sets}
$\{B(x,\varepsilon)\}_{\varepsilon\ge0}$, one per input, indexed by a radius
$\varepsilon\ge0$. We require two structural axioms.

\begin{definition}[Admissible perturbation family]\label{def:pert}
A family $\{B(x,\varepsilon)\}$ is \emph{admissible} if it satisfies
\begin{enumerate}[label=\textup{(P\arabic*)},leftmargin=*]
  \setcounter{enumi}{-1}
  \item \label{P0} $B(x,0)=\{x\}$ for every $x$;
  \item \label{P1} \emph{(nesting)} $\varepsilon_1\le\varepsilon_2$ implies
  $B(x,\varepsilon_1)\subseteq B(x,\varepsilon_2)$.
\end{enumerate}
\end{definition}

Both axioms hold for $\ell_p$ balls $B_p(x,\varepsilon)=\{x':\|x'-x\|_p\le
\varepsilon\}$ and, in the prompt setting, for families of adversarial prompts
measured by suffix length, edit distance, or token-substitution radius; this is the
sense in which our results specialize to the discrete adversarial-prompt threat
models of \cite{Zou2023,Chao2024}.

\subsection{Scores and preservation ratios}
The model is scored on two axes.

\begin{definition}[Scores and worst-case scores]\label{def:scores}
Let $C(f,x)\in[0,1]$ and $A(f,x)\in[0,1]$ be the clean \emph{scientific-correctness}
and \emph{ethical-alignment} scores at $x$; the alignment score aggregates
dimensions such as bias, reproducibility, transparency, integrity, and privacy, and
is of the form $A=1-\mathrm{ASR}$ in the sense of \cite{Chao2024}. The
\emph{worst-case scores} over $B(x,\varepsilon)$ are
\[
  \bar C(f,x,\varepsilon)=\min_{x'\in B(x,\varepsilon)}C(f,x'),
  \qquad
  \bar A(f,x,\varepsilon)=\min_{x'\in B(x,\varepsilon)}A(f,x').
\]
\end{definition}

By \ref{P0} and \ref{P1}, $\bar C\le C$, $\bar A\le A$ pointwise, with equality at
$\varepsilon=0$.

\begin{definition}[Component robustness and composite coefficient]\label{def:rc}
Fix a floor $\tau\in(0,1)$. The \emph{component robustness ratios} are
\[
  \rho_C(f,x,\varepsilon)=\clip\!\Big(\tfrac{\bar C(f,x,\varepsilon)}{\max(C(f,x),\tau)},0,1\Big),
  \qquad
  \rho_A(f,x,\varepsilon)=\clip\!\Big(\tfrac{\bar A(f,x,\varepsilon)}{\max(A(f,x),\tau)},0,1\Big).
\]
For a weight $\lambda\in[0,1]$, the \emph{composite ethical-robustness coefficient}
is
\[
  \RC_\lambda(f,x,\varepsilon)=\lambda\,\rho_C(f,x,\varepsilon)+(1-\lambda)\,\rho_A(f,x,\varepsilon),
\]
and its population version is $\RC_\lambda(f,\varepsilon)=\E_x[\RC_\lambda(f,x,\varepsilon)]$.
\end{definition}

The floor $\tau$ guards against division by a vanishing clean score; it is active
only on the event $\{C<\tau\}$ or $\{A<\tau\}$, and the outer $\clip$ ensures
$\rho_C,\rho_A\in[0,1]$ regardless.

\begin{definition}[Deficit quantities]\label{def:deficit}
For a score $S\in\{C,A\}$ write the \emph{natural deficit}
$e_\nat^S=\E[1-S]$, the \emph{robust deficit} $e_\rob^S=\E[1-\bar S]$, and the
\emph{stability gap} $b^S=\E[S-\bar S]\ge0$.
\end{definition}

\subsection{Regularity and constraint hypotheses}
Several results use clean-competence floors and Lipschitz regularity.

\begin{definition}[Clean floors and Lipschitz constants]\label{def:reg}
We say the model has \emph{clean floors} $c_{\min},a_{\min}\in(0,1]$ if
$C(f,x)\ge c_{\min}$ and $A(f,x)\ge a_{\min}$ almost surely. We say $C(f,\cdot)$ is
$L_C$-Lipschitz on $B(x,\varepsilon)$ if $|C(f,x')-C(f,x'')|\le L_C\,d(x',x'')$ for
all $x',x''\in B(x,\varepsilon)$, and similarly $L_A$ for $A(f,\cdot)$, where
$d(\cdot,\cdot)$ is the metric defining the perturbation set.
\end{definition}

The mechanism connecting ethical training to robustness is modeled after the
constrained-Markov-decision-process (CMDP) feasibility that constrained policy
optimization enforces \cite{Achiam2017}. Concretely, we distinguish two training
regimes through their alignment behavior under the adversary.

\begin{definition}[Constraint and collapse models]\label{def:ec}
The \emph{ethically constrained} model $f_\eth$ satisfies the alignment constraint
\begin{equation}\label{eq:EC}
  \E[\,1-\bar A(f_\eth,\cdot,\varepsilon)\,]\le d \tag{EC}
\end{equation}
for a tolerance $d\in[0,1)$, equivalently $\E[\bar A_\eth]\ge 1-d$. The
\emph{accuracy-only} model $f_\acc$ may collapse under the adversary:
\begin{equation}\label{eq:CO}
  \E[\,\bar A(f_\acc,\cdot,\varepsilon)\,]\le \bar a \tag{CO}
\end{equation}
for a collapse level $\bar a\in[0,1]$.
\end{definition}

Constraint \eqref{eq:EC} is exactly the CMDP feasibility $J_C(\pi_\eth)\le d$ of
\cite[Def.~2, Prop.~2]{Achiam2017}, realizable in practice by the constitution of
Constitutional AI \cite{Bai2022a} or the KL-regularized reward of RLHF
\cite{Ouyang2022,Bai2022b}; the collapse \eqref{eq:CO} models the empirically
documented jailbreak vulnerability of unconstrained models \cite{Zou2023}.

\subsection{Prerequisite results}
We use, without reproof, the following classical statements.

\begin{theorem}[TRADES decomposition {\cite[Def.~2.1--2.3]{Zhang2019}}]\label{thm:trades}
For a binary classifier the robust error decomposes exactly as
$R_\rob(f)=R_\nat(f)+R_\mathrm{bdy}(f)$, a clean-error term plus a non-negative
boundary term.
\end{theorem}

\begin{theorem}[Certified radius {\cite[Thm.~1]{Cohen2019}}]\label{thm:cohen}
Let $g$ be the Gaussian smoothing of a base classifier with variance $\sigma^2$. If
at $x$ the top and runner-up class probabilities under the smoothing noise satisfy
$p_A\ge p_B$, then $g$ is constant on the $\ell_2$-ball of radius
$R=\tfrac{\sigma}{2}\big(\Phi^{-1}(p_A)-\Phi^{-1}(p_B)\big)$ about $x$.
\end{theorem}

\begin{theorem}[Lipschitz margin bound {\cite[Cor.~2.1]{Weng2018}}]\label{thm:clever}
If the class-margin function of a classifier is $L$-Lipschitz on $B_p(x,R)$, then
the minimum distortion required to change its prediction at $x$ is at least the
margin at $x$ divided by $L$.
\end{theorem}
